\documentclass{article}

\usepackage{fullpage}
\usepackage{hyperref}
\usepackage{enumitem}

\begin{document}
\section*{ASEN 5264 Final Project Report Assignment}

The final assignment for the class is report on your final project. The goal is to produce a paper that is similar to a conference paper (but probably with narrower scope and less novelty).

\subsection*{Formatting}

The length must be between 4 and 8 pages \emph{including references}. You may include appendices with additional details, but we may not have time to read them. I recommend using the IEEE format - templates can be found here: \url{https://www.ieee.org/conferences/publishing/templates.html}. You may use other formats, but they should be single spaced and have margins similar to the IEEE format (i.e. versions suitable for final submission to a conference rather than drafts).

\subsection*{Target Audience}

The target audience that you should write to is a peer in this class. You can assume the reader will be familiar with the basic concepts introduced in the class, but you should describe the details of the work you did.

\subsection*{Outline}

I recommend using the following outline (you don't have to use this exactly, but all of the info should be included):

\begin{enumerate}[nosep]
    \item Introduction
    \item Background and Related Work
    \item Problem Formulation (make sure to clearly define S, A, T, and R if you are using an MDP example)
    \item Solution Approach
    \item Results
    \item Conclusion (and possible Future Work if you want to describe potential extensions)
    \item Contributions and Release (see below)
\end{enumerate}

\section*{Rough Rubric (May be adjusted)}

\begin{itemize}
    \item \textbf{Project Design and Scope} (20 pts): Was the project appropriately designed and scoped for the course? Does the project match what was outlined in the proposal?
    \item \textbf{Project Execution} (20 pts): Did the project accomplish the main goals outlined in the proposal? If not, did the project at least accomplish the minimum working example? Were the DMU methods effective in solving the problem?
    \item \textbf{Communication} (50 pts): Does the report clearly communicate the outcome of the project? (Note that some projects may not have the exact structure below)
        \begin{itemize}
            \item \textbf{Introduction, Background, and Related Work} (8 pts): Clearly explain why this is an interesting or important problem and any concepts the target audience would need to understand the rest of the report.
            \item \textbf{Problem Formulation} (12 pts): Describe the problem that you are solving in detail and with minimal ambiguity. In most (but not all) cases, the easiest way to do this is with math. For example, make sure to clearly define $S$, $A$, $T$, and $R$ if you are using an MDP example. If you have a complex problem, it may be easier to define the variables in the state (e.g. ``$s = (x,y,\theta)$ where $x$ is \ldots''), and it may be easier to define a generative model or episodic simulator (e.g. ``The state transitions are implicitly defined as a simulator which integrates the following equations\ldots''). Figures can help.
            \item \textbf{Solution Approach} (12 pts): Describe the solution approach in detail and with minimal ambiguity. In most (but not all) cases, the easiest way to do this is with pseudocode and/or simple figures like the DQN description from class. Make sure to describe implementation choices such as rollout policies or network architectures.
            \item \textbf{Results} (10 pts): Describe the results of your project with plots and tables. Quantify uncertainty if at all possible. Present the most important results in the most prominent places and highlight the key takeaways (e.g. by bolding the best performers in a table).
            \item \textbf{Discussion and Conclusions} (8 pts): What are the takeaways from the results? Did your solution perform well? Why or why not, and what would you do differently next time?
        \end{itemize}
    \item \textbf{Required Items} (10 pts): See below.
\end{itemize}

\subsection*{Required Items}
\begin{itemize}
    \item You must conduct a basic familiarization of literature related to your project. Your report should contain \textbf{at least 5 citations of related published work} and a brief description of how these relate to your project in the introduction or background section.
    \item At the end of the report include a few sentences about the \textbf{contributions of each team member} (e.g. ``Bob collected and cleaned the data, Alice implemented the algorithm", etc.)
    \item At the end of the report, \textbf{indicate whether you are willing to share the report publicly} for other students to benefit from. I highly recommend choosing to share since future students may be able to build off of or be inspired by your work. \textbf{Please include one of the following release sentences}: ``The authors grant permission for this report to be posted publicly." or "The authors do NOT grant permission for this report to be posted publicly.''
    \item In a brief paragraph or bullets at the end of the report, clearly \textbf{indicate which algorithms you implemented yourself} vs used off-the-shelf. ``Implement'' means that you wrote the core logic of the algorithm yourself rather than calling an existing library or package. Using libraries for supporting functionality (e.g. linear algebra, data structures, neural networks) is included, but the main algorithmic logic should be your own if you claim that you implemented it.
    \item If you are in a group, you must \textbf{submit as a group on gradescope}.
\end{itemize}

\section*{AI (Large Language Model) Policy}

The AI policy for this assignment is the same as for the course:

I approach the class with the assumption that all students have a strong desire to learn rather than just to complete the assignments. In this course, you are highly-encouraged to use AI (specifically large language models such as Claude) as a \emph{learning tool}. You can ask it questions or have it assist in writing some code, but any prose or mathematical work that you turn in must be written by you, and any code must be written by or closely directed by you. You must understand everything that you turn in.

In particular, for this assignment, it is important that the report is written by you. You are encouraged to consult AI to get ideas and ask it for help, but you are prohibited from directing it to write the report.

\end{document}
